\section{Introduction}
\label{sec:intro}

Large language models are increasingly evaluated on knowledge-intensive
benchmarks, but the most widely-used resources --- MMLU
\citep{hendrycks2021mmlu}, GPQA \citep{rein2024gpqa}, and their
successors \citep{srivastava2023bigbench, zhong2024agieval, phan2025hle}
--- share three weaknesses: they aggregate breadth at the expense of
\emph{depth} within any single discipline; they are typically authored
by a small team using one generation pipeline, exposing systematic
blind-spots and stylistic regularities that frontier models can learn
to exploit \citep{panickssery2024selfpref, zheng2023judging}; and the
questions are increasingly suspected of contamination through web-scale
pre-training corpora \citep{sainz2023contamination,
magar2022contamination, dodge2021documenting}. As models grow more
capable, the marginal value of a benchmark depends increasingly on
\emph{how} it was constructed.

We argue that domain-specialised benchmarks built around bodies of
knowledge that are \emph{externally validated by humans} --- through
certification, regulation, or peer-reviewed practice --- offer a useful
complement to broad knowledge benchmarks. Examples already exist in
medicine \citep{jin2021medqa, jin2019pubmedqa}, law
\citep{guha2023legalbench}, and finance \citep{islam2023financebench}. We
extend this line of work to a domain that has so far received little
attention from the benchmarking community, despite being unusually
well-defined: \textbf{wine}.

\paragraph{Why wine.} Wine integrates plant biology, geology and
climate science, organic chemistry, food microbiology, sensory
neuroscience, regulatory law, and business economics. It is one of
very few domains for which inter-governmental and governmental sources
publish structured taxonomic data (INAO, OIV, TTB, regional consortia),
university research groups (UC Davis, Geisenheim, Bordeaux Sciences
Agro) produce open scholarly literature, and professional certification
bodies (WSET Levels 1--4, Court of Master Sommeliers, Institute of
Masters of Wine) define a graded competence ladder. These properties
make it possible to build a benchmark whose every question is
traceable to a verifiable source while still spanning the
recall--reasoning--application spectrum that distinguishes a strong
taster, winemaker, or buyer from a memorised glossary.

\paragraph{Contributions.}
\begin{enumerate}\itemsep -1pt
  \item \textbf{A 3{,}266-question wine benchmark} across six pillars
        and four difficulty tiers, built from 38{,}104 atomic
        provenance-tagged facts (Section~\ref{sec:dataset}).
  \item \textbf{An LLM-driven, fact-grounded generation pipeline} with
        five strategies $\times$ five generator models, per-strategy
        quotas, and an explicit closed-book pre-screen --- LLMs
        reformat verified facts, never supply them.
  \item \textbf{A nine-agent automated audit} (Section~\ref{sec:audit})
        calibrated against a human gold sheet via Cohen's $\kappa$;
        the audit removed 341 questions and re-labelled difficulty on
        1{,}259, and surfaced one agent (closed-book solvability) that
        over-reports leakage relative to humans, which we keep with
        explicit disclosure rather than drop.
  \item \textbf{A bias-aware evaluation across sixteen frontier
        configurations} (Section~\ref{sec:eval}) reporting overall
        accuracy, reasoning-mode lift, a Self-Preference Score,
        cost-efficiency Pareto, and a closed-book vs.\ source-grounded
        contrast that isolates parametric recall from contextual
        reasoning.
\end{enumerate}

\noindent OenoBench --- corpus, audit findings, construction code, and
human-review app --- ships under CC-BY-SA-4.0; release URLs,
reproducibility scripts, and the per-source licensing table are
collected in Appendix~\ref{app:artifacts}.
